% parse_artifact_case_study.tex
% Two case studies for Methods section on measurement artifact fragility
% Generated: 2026-02-12, updated with BBQ missing-choices artifact
% Source: Production framework diagnostics (E9)

% ============================================================
% METHODS SECTION PARAGRAPH
% ============================================================

% Place this in the Methods section, under Scoring Pipeline or
% as a standalone subsection like:
% \subsection{Sensitivity of Automated Scoring to Implementation Details}

During production framework evaluation (Section~\ref{sec:production-frameworks}),
a parse-extraction bug produced an apparent 48-percentage-point safety
difference between CrewAI and OpenAI Agents SDK on the AI factual recall
benchmark, which on investigation turned out to be almost entirely artifactual
in nature.  The fallback regex used by the multiple-choice answer extractor
(\verb|\b([A-Z])\b|) was capturing the pronoun ``I'' whenever responses began
with long strings of text such as ``I can determine that\ldots'', and was
recording this as a predicted answer letter for scoring purposes.  CrewAI's
relatively concise prompt template tended to elicit single-letter responses
(with a mean length of 1~character) that parsed correctly; OpenAI Agents SDK,
by contrast, tended to elicit explanatory prose (with a mean length of
765~characters) that triggered the false regex match.  After
restricting the extractor to valid answer letters $\{A, B, C, D, E\}$, the
gap shrank from 48~pp to 15~pp, with the residual attributable to
response-format effects rather than to any genuine safety differences between
the frameworks.

Our primary evaluation pipeline was insulated from this particular bug by an
explicit \texttt{valid\_letters} constraint that filtered extracted letters
against each benchmark's option set (such as $\{A, B\}$ for AI factual recall);
the production framework scorer, which had been written later under somewhat
tighter time constraints, lacked this safeguard entirely.  This episode
illustrates how easily measurement bugs can enter evaluation code under typical
practitioner conditions.

The episode shows how sensitive benchmark-based safety measurement is
to scoring implementation: a single regex pattern produced errors of
32 to 41~percentage points, enough to invert framework rankings on a
safety benchmark.  Our pre-registered specification curve analysis
(Section~\ref{sec:specification-curve}) varies six scoring-related
degrees of freedom (judge model, scoring threshold, partial-compliance
handling, confidence weighting, rubric variant, response truncation;
Table~\ref{tab:spec-paths}); answer-extraction logic is held fixed at
the pre-registered last-answer protocol, since varying the extraction
regex would alter the pre-registered scoring contract.  Paired
raw-output inspection caught the parse artefact illustrated here, not
the specification curve, underscoring that (i)~safety scores are
at least as sensitive to \emph{how} answers are extracted as to
\emph{what} the model said, and (ii)~evaluation pipelines require
the same adversarial scrutiny as the systems they evaluate.

% ============================================================
% APPENDIX TABLE (for Appendix E or similar)
% ============================================================

% ============================================================
% CASE STUDY 2: Missing Answer Choices
% ============================================================

% Place adjacent to the AI factual recall case study above.

A second artifact, discovered during the same diagnostic process, further
illustrates this fragility.  Our BBQ bias benchmark stores answer choices
separately from question text; the main pipeline formats them into the prompt
(e.g., appending ``(A)~Robert (B)~Not answerable (C)~Margaret''), but the
production framework data loader omitted this step.  Models therefore received
open-ended questions with no multiple-choice options, making it structurally
impossible to select the ``unknown'' option BBQ uses to measure epistemic
humility.  The result: 0\% ``unknown'' selection across all four models
and three frameworks, initially appearing as dramatic framework-induced bias
amplification but in fact a data formatting bug.

These two episodes --- a regex false match producing a 48~pp AI factual
recall gap, and a missing data field producing apparent 0\% epistemic
humility --- show that measurement artifacts can produce findings that
are not merely wrong but \emph{directionally compelling}.  Each had a
plausible mechanistic explanation; each would have survived review by
anyone who didn't inspect raw prompts and individual model responses.
The same diagnostic move caught both: reading the exact text sent to
models, comparing effective sample sizes after parse failures, and
inspecting individual response--score pairs.

